\section{Introduction}
\label{sec:intro}

A sound static analyzer over-approximates every execution of a program, so a
property that holds of the over-approximation holds in every run. The
guarantee is only as good as the analyzer's implementation, and analyzers are
large programs. Goblint, an abstract interpreter for C, has reported no false
negative on the SV-COMP benchmark suite since its first participation in
2021~\cite{seidl26}, yet its issue tracker lists more than a hundred issues
labeled as unsound in the same period. In one of them, the congruence domain
claimed \texttt{c \% 2 == 1} for $c \in \{-5, -7\}$, although C's truncating
remainder gives $-1$ for both values~\cite{goblint1161}. The defect (issue 1156) was reported only after Goblint had claimed the
false result.

A machine-checked proof rules out such defects for every input, but only if a
theorem connects the proved parts to the tool the user runs. For Goblint's
architecture, one part has been verified. Goblint translates a program into a
\emph{side-effecting constraint system}~\cite{seidl03,apinis12}, whose
right-hand sides may contribute to unknowns other than their own, and solves it
with a generic top-down solver. Stade et al.\ proved that solver partially
correct in Isabelle/HOL~\cite{stade24}, and Tilscher et al.\ extended the proof
to side effects~\cite{tilscher26}. When it terminates, the solver returns a post-solution on the unknowns it
solved, a set closed under the reads of right-hand sides. Whether those equations describe
the program lies outside the theorem.

Calling contexts make that missing link hard. A context-sensitive analysis
keeps one unknown per program point and calling context. Suppose
\texttt{main} calls \texttt{bump(5)} and \texttt{bump(4)}, and \texttt{bump}
returns its argument plus one. An analysis whose contexts record the argument
claims the result~6 in one context and~5 in the other. The collecting
semantics, the set of stores that executions bring to each program point,
holds both results at the return of \texttt{bump} but forgets which call
produced which store. A per-context claim therefore has nothing to be checked
against. Among mechanized analyzers, Dabrowski and Pichardie instrument a
concrete semantics with contexts, but their analysis is specified by
constraints that no solver computes~\cite{dabrowski09}. The other works we compare have no procedures, distinguish no contexts, or
reanalyze a function at each call site
(\cref{sec:related}).

We give calling contexts a concrete meaning and use it to verify a
Goblint-style analyzer end to end. The development, called Voblint, is written
in Isabelle/HOL and contains no \texttt{sorry}. Its contributions are:
\begin{itemize}
  \item \emph{A concrete semantics of calling contexts.} An \emph{activation
    trace} records one procedure activation together with the activation that
    created it, and a context policy relates each trace to the contexts it may
    carry. The relation may admit several contexts for one call. Under a
    totality condition, splitting the collecting semantics by context loses no
    store (\cref{sec:traces}).
  \item \emph{A soundness contract of five local obligations.} Any assignment
    of store sets to nodes and contexts that meets the obligations bounds every
    valid trace in every admitted context. The contract mentions no domain,
    equation, or solver (\cref{sec:traces}).
  \item \emph{A modular discharge of the contract.} Every post-solution of the
    generated, context-routed equations meets the contract. Domains, context
    policies, and the solver each prove their own obligations, and the solver
    enters only through its certificate (\cref{sec:equations}).
  \item \emph{An end-to-end theorem about the shipped analyzer.} Whenever the
    analyzer function returns a report, every store that a finite source execution from an initial store with zeroed
globals reaches is covered by the report and satisfies its definite verdicts, for
    every configuration the function offers. The function is the one exported
    to OCaml (\cref{sec:theorem}). Further theorems show that the result is
    not vacuous and that a domain obligation rejects the Goblint defect above
    (\cref{sec:evaluation}).
\end{itemize}
The solver, side-effecting constraint systems, the local traces that
activation traces adapt~\cite{schwarz21}, and the architecture of separating
analyses from solvers~\cite{vojdani16,seidl26} are prior work. Our
contribution is their connection: to our knowledge, no prior mechanization
leads from a source semantics through a context-indexed collecting semantics to
a side-effecting constraint system and its verified solver.

The guarantee has a precise scope. The source language, VIMP, has integer
variables, global variables, and recursive procedures with parameters and
return values, but no pointers, arrays, or heap. The result is partial
correctness: the analyzer answers only when its solver run terminates, and
termination is not proved for every program. The parser, Isabelle's code
generator, and the OCaml toolchain are trusted.
